\section{Теоретические основы}
\label{sec:theory}

В обзоре литературы было выделено пять взаимосвязанных направлений:
репликация и коды стираний, локальное восстановление, температурные политики
хранения, гибридные схемы и конвертируемые коды. Для предлагаемого решения эти
направления нужно объединить в одну модель. В данной главе вводятся
обозначения, которыми далее описываются состояния схем хранения, чтение при
отказах, температура данных и критерий выбора перехода между схемами.

\subsection{Модель распределённого хранения и отказов}

Рассматривается распределённое хранилище, состоящее из множества узлов. Каждый
узел хранит физические блоки данных и имеет ограниченный объём дискового
пространства. В рассматриваемой модели один узел считается одним доменом отказа: если
узел недоступен, все размещённые на нём блоки недоступны одновременно.
Такое допущение не описывает все возможные причины отказов, но удобно для
анализа размещения и восстановления, поскольку именно узел является минимальной
единицей, на которой в системе принимается решение о размещении блока.

В работе используется следующая терминология. Объект разбивается на
последовательность логических блоков одинакового размера. Логический блок
соответствует единице чтения, для которой рассчитывается ожидаемая стоимость
чтения \(E[io]\). Логические блоки раскладываются на физические блоки, которые
являются минимальными единицами размещения на узлах. Помимо блоков данных в
схеме могут присутствовать проверочные блоки, которые используются для
восстановления при отказах. Группа логических блоков образует полосу
кодирования. Для каждой
полосы хранилище поддерживает схему избыточности: набор реплик,
локальных проверочных блоков и глобальных проверочных блоков. В
постановке с дозаписью без изменения уже сохранённых данных, принятой в данной
работе, пользовательские данные после записи не изменяются. Поэтому основными
операциями становятся чтение существующих данных и фоновое изменение схемы
хранения.

Для анализа отказов используется сглаженная оценка доли недоступных доменов
отказа. Если в кластере \(H\) доменов отказа и в момент \(t\) недоступно
\(K(t)\) из них, наблюдаемая доля отказов равна
\[
    q_{obs}(t)=\frac{K(t)}{H}.
\]
Чтобы кратковременные изменения состояния кластера не приводили к резким
переходам схем туда и обратно, используется экспоненциальное сглаживание:
\[
    \hat q(t)=
    \alpha \hat q(t_0) + (1-\alpha)q_{obs}(t),
    \qquad
    \alpha = 2^{-\frac{t-t_0}{T_q}},
\]
где \(T_q\) --- период полураспада оценки отказов. При первом измерении
полагается \(\hat q(t)=q_{obs}(t)\).

Если схема занимает \(n_S\) независимых доменов отказа, а
\(\hat K=\operatorname{round}(\hat q H)\) --- сглаженная оценка числа
недоступных доменов в кластере, вероятность ровно \(f\) отказов среди доменов
этой схемы можно оценить гипергеометрически:
\[
    p_f(S)=
    \frac{
        \binom{\hat K}{f}\binom{H-\hat K}{n_S-f}
    }{
        \binom{H}{n_S}
    }.
\]
Если схема занимает малую долю кластера, допустима более простая биномиальная
аппроксимация \(p_f=\binom{n_S}{f}\hat q^f(1-\hat q)^{n_S-f}\). Далее
предполагается, что блоки, совместно нужные для восстановления, размещаются
в разных доменах отказа.
Так как в работе сравниваются схемы, рассчитанные на три отказа, при
практической оценке можно использовать усечённое распределение:
\[
    p_0^*=p_0,\qquad
    p_1^*=p_1,\qquad
    p_2^*=p_2,\qquad
    p_{3+}=1-p_0-p_1-p_2.
\]
Последний член означает вероятность трёх или более отказов и задаёт
консервативную оценку стоимости чтения в наиболее тяжёлом учитываемом случае.

\subsection{Репликация и глобальная избыточность}

Репликация является наиболее простой формой избыточности. Если каждый
логический блок хранится в нескольких репликах, чтение обычно требует одной
операции: достаточно выбрать любую доступную реплику. При отказе одной реплики
данные можно прочитать из другой реплики и восстановить утраченную. Главный
недостаток репликации состоит в расходе места: трёхкратная репликация требует
примерно втрое больше дискового пространства, чем исходный объём данных.

Коды стираний уменьшают этот расход. В классической схеме Рида--Соломона
данные рассматриваются как набор символов над конечным полем, а проверочные
символы вычисляются как линейные комбинации исходных символов
\cite{reed1960polynomial}. Для систем хранения такие коды часто выступают
базовой точкой сравнения: \enquote{Reed-Solomon codes are the standard design
choice}~\cite{sathiamoorthy2013xoring}. Если схема содержит \(d\) блоков данных и
\(g\) глобальных проверочных блоков, то при свойстве максимального расстояния
по столбцам (свойстве разделимости с максимально достижимым расстоянием) можно
восстановить данные после потери до \(g\) произвольных блоков. Практические
системы используют такие коды для снижения расхода места, но платят за это
более дорогим восстановлением: при недоступности нужного блока приходится
читать несколько других блоков схемы
\cite{huang2012azure,shen2024erasurecoding}.

В дальнейшем глобальными проверочными блоками будем называть проверочные
блоки, построенные по всем \(d\) блокам данных. Они обеспечивают
устойчивость к нескольким отказам, но восстановление через них обычно требует
обращения к широкому набору блоков.

\subsection{Локально восстанавливаемые коды}

Локально восстанавливаемые коды уменьшают стоимость восстановления отдельных
блоков. Их идея состоит в том, что блоки данных разбиваются на
локальные группы, а для каждой группы хранится отдельный проверочный блок.
Если в такой группе потерян один блок данных, его можно восстановить,
прочитав только остальные блоки этой группы и локальный проверочный
блок. Теоретическая модель локальности для кодовых символов рассмотрена в
\cite{gopalan2012lrc}, а практическая мотивация для систем хранения описана в
\cite{huang2012azure,sathiamoorthy2013xoring}. В работах о LRC это свойство
формулируется так: \enquote{can repair any codeword symbol using a
small number of other symbols}~\cite{kong2024lrcconvertible}.

Пусть полоса содержит \(d\) блоков данных и \(l\) локальных проверочных
блоков. Если \(l>0\) и \(d\) кратно \(l\), размер локальной группы равен
\[
    s = \frac{d}{l}.
\]
Один локальный проверочный блок покрывает одну группу из \(s\) блоков
данных. При одиночном отказе внутри группы стоимость восстановления становится
порядка \(s\), а не \(d\). Если же отказов несколько и они затрагивают
несколько блоков одной группы, локального проверочного блока может быть недостаточно,
и тогда требуется восстановление через глобальные проверочные блоки.

Таким образом, локальность задаёт ещё одну ось компромисса. Увеличение числа
локальных групп снижает стоимость типичных одиночных восстановлений, но
добавляет проверочные блоки и усложняет размещение. Для данной работы
важно, что локальные проверочные блоки можно рассматривать не только как
часть фиксированного кода, но и как промежуточное состояние между
репликацией и более компактным кодовым хранением.

\subsection{Конвертируемые коды и LRCC}

Если схема хранения меняется после записи данных, возникает задача
перекодирования. Наивный переход требует прочитать исходные данные, заново
вычислить все проверочные блоки и удалить старую избыточность.
Конвертируемые коды формализуют эту задачу: они задают начальный и конечный
коды так, чтобы переход между ними требовал как можно меньшего числа чтений и
записей \cite{maturana2022convertible}. Их цель состоит в том, чтобы выполнять
переходы с \enquote{significantly less resources than the default approach of
re-encoding}~\cite{maturana2022convertible}. Поэтому важна не только устойчивость
каждой отдельной схемы, но и структура перехода между схемами.

Локально восстанавливаемые конвертируемые коды, или LRCC, появляются как
сочетание двух идей: начальная и конечная схемы являются локально
восстанавливаемыми, а переход между ними рассматривается в модели
конвертируемых кодов \cite{kong2024lrcconvertible}. Поэтому LRCC логически
удобнее вводить именно после LRC и конвертируемых кодов. Это не просто ещё
одна разновидность локально восстанавливаемого кода, а LRC-схема, для которой
существенна возможность эффективного перехода из одного набора параметров в
другой.

Morph использует эту идею как системный строительный блок
\cite{kim2024morph}. В большинстве файловых систем такая операция не является
нативной: \enquote{transcode is not directly supported as a native DFS
operation}~\cite{kim2024morph}. Важное наблюдение Morph состоит в том, что проверочные
блоки, записанные на ранней стадии жизни файла, можно выбирать так, чтобы
позже они становились частью более холодной схемы. Например, проверочный
блок ранней схемы может быть переиспользован как локальный проверочный
блок конечной LRC/LRCC-схемы, а остальные проверочные блоки могут быть
объединены в глобальные. Для данной работы Morph важен не как источник новой
теории LRCC, а как пример того, что жизненный цикл данных можно проектировать
с учётом будущих переходов.

\subsection{Гибридная схема Hy(r, LRCC(d,l,g))}
\label{sec:theory-hybrid}

Для описания состояний в данной работе используется гибридная схема
\[
    Hy(r, LRCC(d,l,g)).
\]
Мотивация такого объединения близка к HyRES: даже когда репликация применяется
для горячих данных, а коды стираний --- для холодных, эти подходы часто
рассматриваются раздельно: \enquote{replication and erasure coding techniques
... each is designed in isolation}~\cite{lucani2025hyres}. В данной работе они, напротив, описываются как части
одного пространства состояний.
Здесь \(d\) --- число блоков данных в полосе, \(l\) --- число локальных
проверочных блоков, \(g\) --- число глобальных проверочных блоков, а
\(r\) --- число дополнительных реплик каждого блока данных поверх
LRCC-представления. Сама LRCC-часть всегда содержит \(d\) блоков данных,
поэтому число доступных прямых размещений каждого логического блока равно
\(r+1\): один блок находится в LRCC-представлении, ещё \(r\) размещаются как
реплики.

Такое обозначение позволяет описывать репликацию и кодирование в одном
пространстве состояний. Изменение \(r\) отвечает за обмен реплик на
проверочные блоки и обратно, изменение \(g\) --- за число глобальных
проверочных уравнений, а изменение \(l\) --- за число локальных групп. Если
\(l>0\), требуется, чтобы \(d\) было кратно \(l\), а размер локальной группы
\(d/l\) оставался в допустимых пределах. Это исключает вырожденные схемы и
ограничивает число возможных состояний, которые должен рассматривать
планировщик.

Помимо изменения параметров одной полосы, в этом пространстве естественно
описываются операции разделения и объединения схем. Разделение уменьшает ширину
полосы и может быть выгодно при нагреве разных частей объекта. Объединение
совместимых схем одного объекта, наоборот, увеличивает ширину и помогает
снизить долю избыточности для холодных данных. Для таких переходов недостаточно
совпадения параметров \(r,d,l,g\): требуется совместимость размещения,
локальных групп и целевого состояния.

Для семейства схем, рассчитанного на сохранение читаемости при любых трёх
отказах в абстрактной модели, используется параметрическое условие
\[
    r + g + \mathbf{1}_{l>0} = 3,
\]
где \(\mathbf{1}_{l>0}\) равна единице, если в схеме есть локальный
проверочный блок. Интуитивно \(r\) реплик защищают чтение от потери отдельных прямых
размещений, \(g\) глобальных проверочных блоков дают независимые уравнения для
восстановления, а наличие локального проверочного блока помогает при одиночном отказе
внутри локальной группы. Это условие не является самостоятельной теоремой о
корректности любой схемы с такими параметрами. Оно задаёт выбранный уровень
избыточности, который должен сопровождаться корректным размещением,
разнесением совместно необходимых блоков по разным доменам отказа и
независимыми глобальными проверками.

Например, одной из возможных траекторий охлаждения данных является постепенный
обмен реплик на проверочные блоки: сначала добавляется локальный проверочный
блок, затем глобальные проверочные блоки, а лишние реплики удаляются
только после появления
заменяющей избыточности. Эта траектория является частным случаем общего
пространства переходов, а не единственной допустимой цепочкой жизненного цикла.

\subsection{Чтение при отказах и ожидаемая стоимость чтения}
\label{sec:theory-read-cost}

В нормальном случае чтение одного логического блока требует одной
физической операции чтения: система обращается к доступной реплике или к
соответствующему блоку данных в LRCC-представлении. При отказах возможны
две разные операции. Чтение с восстановлением --- это синхронное восстановление
данных на пути пользовательского чтения без изменения размещения. Фоновое
восстановление --- это восстановление утраченного физического блока на
новом узле с обновлением метаданных. В данной работе эти операции
разделяются: чтение с восстановлением нужно для корректного обслуживания
чтения, а фоновое восстановление относится к эксплуатации физического состояния
хранилища и не является переходом между схемами избыточности.

Для выбора схемы планировщику нужна ожидаемая стоимость чтения. Обозначим её
как \(E[io]\). Это математическое ожидание числа физических чтений, необходимых
для чтения одного логического блока, с учётом вероятностей отказов:
\[
    E[io \mid S] =
    p_0^* io_0(S) + p_1^* io_1(S) + p_2^* io_2(S) + p_{3+} io_3(S),
\]
где \(S\) --- состояние схемы, \(p_0^*,p_1^*,p_2^*,p_{3+}\) --- вероятности
после описанного усечения распределения, а \(io_f(S)\) --- средняя стоимость
чтения при условии, что внутри схемы произошло ровно \(f\) отказов. Величины
\(io_f\) сами по себе не являются итоговым \(E[io]\): итоговое ожидание
получается только после взвешивания на вероятности \(p_0^*,p_1^*,p_2^*,p_{3+}\).

Для гибридной схемы положим
\[
    c = r + 1,\qquad N = d + r + l + g.
\]
Здесь \(c\) --- число прямых размещений читаемого символа данных: один
LRCC-блок данных и \(r\) реплик. Величина \(N\) --- число доменов отказа,
занятых схемой при размещении. Принято допущение, что каждая реплика всей
полосы хранится целиком на одном узле в виде последовательности блоков, а не
разбрасывается на \(d\) независимых узлов. Поэтому реплики увеличивают число
прямых размещений каждого символа данных, но добавляют к \(N\) только
\(r\) доменов отказа. При выборе другой политики размещения реплик ---
например, по одной копии каждого блока данных на отдельном узле --- формулы для
\(N\) и условные стоимости \(io_f\) ниже потребуют пересмотра.
Если локальных групп нет (\(l=0\)), то восстановление при потере всех прямых
размещений читаемого символа происходит через глобальную часть схемы. Тогда
\[
    io_f =
    \frac{
        \left(\binom{N}{f} - \binom{N-c}{f-c}\right)
        + d \binom{N-c}{f-c}
    }{
        \binom{N}{f}
    },
\]
где биномиальные коэффициенты с некорректными параметрами считаются равными
нулю. Первая часть числителя соответствует случаям, когда хотя бы одно прямое
размещение читаемого символа доступно, а вторая --- случаям, когда все прямые
размещения потеряны и требуется восстановление через \(d\) блоков.

Если \(l>0\), то появляется локальное восстановление. Пусть
\[
    s = \frac{d}{l}
\]
--- размер локальной группы. Для фиксированного читаемого символа число
паттернов отказов, при которых все его прямые размещения потеряны, но локальное
восстановление всё ещё возможно, обозначим \(B_f\):
\[
    B_f =
    \binom{N-c-s}{f-c}.
\]
После потери \(r\) узлов с репликами и LRCC-блока читаемого символа
локальное восстановление возможно только если доступны локальный проверочный
блок этой
группы и \(s-1\) остальных LRCC-блоков данных в группе. Тогда средняя
стоимость чтения при \(f\) отказах равна
\[
    io_f =
    \frac{
        A_f + s B_f + d G_f
    }{
        \binom{N}{f}
    },
\]
где
\[
    A_f = \binom{N}{f} - \binom{N-c}{f-c},
    \qquad
    G_f = \binom{N-c}{f-c} - B_f.
\]
Здесь \(A_f\) соответствует чтению доступного прямого размещения, \(B_f\) ---
локальному восстановлению стоимостью \(s\), а \(G_f\) --- более дорогому
восстановлению через глобальную часть схемы стоимостью порядка \(d\).

Для наглядности в табл.~\ref{tab:io-fixed-failures} приведены значения
условной стоимости \(io_f\) для \(f=0,1,2,3\). Эти числа показывают, почему
широкие и сильно закодированные схемы выгодны по месту, но хуже при чтении с
восстановлением, а локальное разбиение уменьшает стоимость типичного
восстановления.

\begin{center}
\begin{longtable}{|c|c|c|c|c|c|}
\caption{\label{tab:io-fixed-failures}
Условное число физических чтений при фиксированном числе отказов внутри схемы.}\\
\hline
Схема & \(N\) & \(io_0\) & \(io_1\) & \(io_2\) & \(io_3\) \\
\hline
\endfirsthead
\hline
Схема & \(N\) & \(io_0\) & \(io_1\) & \(io_2\) & \(io_3\) \\
\hline
\endhead
\(Hy(3,LRCC(6,0,0))\) & 9 & 1{,}000 & 1{,}000 & 1{,}000 & 1{,}000 \\
\(Hy(3,LRCC(12,0,0))\) & 15 & 1{,}000 & 1{,}000 & 1{,}000 & 1{,}000 \\
\(Hy(3,LRCC(24,0,0))\) & 27 & 1{,}000 & 1{,}000 & 1{,}000 & 1{,}000 \\
\(Hy(2,LRCC(6,1,0))\) & 9 & 1{,}000 & 1{,}000 & 1{,}000 & 1{,}060 \\
\(Hy(2,LRCC(12,2,0))\) & 16 & 1{,}000 & 1{,}000 & 1{,}000 & 1{,}009 \\
\(Hy(2,LRCC(12,1,0))\) & 15 & 1{,}000 & 1{,}000 & 1{,}000 & 1{,}024 \\
\(Hy(2,LRCC(24,4,0))\) & 30 & 1{,}000 & 1{,}000 & 1{,}000 & 1{,}001 \\
\(Hy(2,LRCC(24,2,0))\) & 28 & 1{,}000 & 1{,}000 & 1{,}000 & 1{,}003 \\
\(Hy(2,LRCC(24,1,0))\) & 27 & 1{,}000 & 1{,}000 & 1{,}000 & 1{,}008 \\
\(Hy(1,LRCC(6,1,1))\) & 9 & 1{,}000 & 1{,}000 & 1{,}139 & 1{,}417 \\
\(Hy(1,LRCC(12,2,1))\) & 16 & 1{,}000 & 1{,}000 & 1{,}042 & 1{,}189 \\
\(Hy(1,LRCC(12,1,1))\) & 15 & 1{,}000 & 1{,}000 & 1{,}105 & 1{,}314 \\
\(Hy(1,LRCC(24,4,1))\) & 30 & 1{,}000 & 1{,}000 & 1{,}011 & 1{,}061 \\
\(Hy(1,LRCC(24,2,1))\) & 28 & 1{,}000 & 1{,}000 & 1{,}029 & 1{,}131 \\
\(Hy(1,LRCC(24,1,1))\) & 27 & 1{,}000 & 1{,}000 & 1{,}066 & 1{,}197 \\
\(Hy(0,LRCC(6,1,2))\) & 9 & 1{,}000 & 1{,}556 & 2{,}111 & 2{,}667 \\
\(Hy(0,LRCC(12,2,2))\) & 16 & 1{,}000 & 1{,}312 & 1{,}925 & 2{,}677 \\
\(Hy(0,LRCC(12,1,2))\) & 15 & 1{,}000 & 1{,}733 & 2{,}467 & 3{,}200 \\
\(Hy(0,LRCC(24,4,2))\) & 30 & 1{,}000 & 1{,}167 & 1{,}582 & 2{,}178 \\
\(Hy(0,LRCC(24,2,2))\) & 28 & 1{,}000 & 1{,}393 & 2{,}167 & 3{,}080 \\
\(Hy(0,LRCC(24,1,2))\) & 27 & 1{,}000 & 1{,}852 & 2{,}704 & 3{,}556 \\
\hline

\end{longtable}
\end{center}

После выбора вероятностей отказов итоговое \(E[io]\) получается взвешиванием
значений \(io_f\). Например, для кластера из \(H=30\) узлов и \(K=3\)
недоступных узлов вероятности \(p_0^*,p_1^*,p_2^*,p_{3+}\) вычисляются
гипергеометрически.

Эта модель намеренно остаётся скалярной. Она оценивает не задержку и не
сетевой трафик напрямую, а ожидаемое число физических чтений. Более детальные
показатели, такие как число задействованных узлов или объём переданных байтов,
могут использоваться как диагностические метрики, но не входят в основную
функцию выбора перехода.

\subsection{Температура данных и функция полезности перехода}
\label{sec:theory-temperature}

Стоимость чтения важна только для тех схем, к которым действительно обращаются.
Поэтому каждой схеме сопоставляется температура \(\lambda_i\), отражающая
текущую интенсивность чтения. Такая постановка опирается на наблюдение, что
реальные нагрузки часто неравномерны: \enquote{OLTP workloads typically
exhibit skewed access patterns}~\cite{levandoski2013hotcold}. В рассматриваемой модели это метрика уровня
схемы: чтение полностью или частично пересекает схему и создаёт вклад,
пропорциональный прочитанной доле \(w\) её логического диапазона. Путь чтения
накапливает такие вклады в счётчиках, а ответственный планировщик периодически
обновляет оценку температуры. Используется экспоненциально взвешенное скользящее
среднее (далее EWMA): между обновлениями применяется отложенное
экспоненциальное затухание. Для оценки горячих и холодных данных такой подход
поддерживается экспериментальным наблюдением, что \enquote{exponential smoothing
produces very accurate estimates}~\cite{levandoski2013hotcold}.
\[
    \lambda_i(t) =
    \lambda_i(t_0) \cdot 2^{-\frac{t-t_0}{T_{1/2}}} + w,
\]
где \(T_{1/2}\) --- период полураспада температуры. В целевой архитектуре
температура остаётся входной метрикой планировщика и может обновляться
асинхронно по накопленным счётчикам чтения. Такая модель не требует
периодически обновлять все схемы: значение пересчитывается только при чтении
или при оценке планировщиком.

Температура трактуется как суммарная интенсивность чтений схемы. Поэтому при
разделении схемы без информации о распределении температуры внутри схемы нагрузка
распределяется равномерно: две дочерние схемы получают по \(\lambda_i/2\). При
объединении двух схем температура новой схемы равна сумме температур
объединяемых схем. Так сохраняется суммарный вклад пользовательской нагрузки в
\(IO_{sum}\).

Для всего кластера ожидаемая нагрузка чтения задаётся суммой
\[
    IO_{sum} = \sum_i \lambda_i E_i,
\]
где \(E_i = E[io \mid S_i]\) --- ожидаемая стоимость чтения текущего состояния
схемы \(i\). Эта величина отличается от средней нагрузки на одну схему:
планировщик должен учитывать именно пользовательскую нагрузку, поэтому горячая
схема с большим \(\lambda_i\) должна влиять на решение сильнее, чем множество
холодных схем.

Для нормировки IO-давления используются глобальные опорные точки, не зависящие
от текущего состояния отдельных схем. Пусть \(T_{total} = \sum_i \lambda_i\) ---
суммарная температура кластера. Тогда нижняя граница IO соответствует хранению
всех данных в виде реплик, при котором \(E[io] = 1\):
\[
    IO_{min} = T_{total} \cdot 1{,}0.
\]
Верхняя граница соответствует самой широкой плотной схеме:
\[
    S_{wide}=\mathrm{LRCC}(d_{max}, 1, 2).
\]
Такая схема занимает наибольшее число доменов отказа среди компактных
состояний и поэтому сильнее всего страдает от отказов в гипергеометрической
модели:
\[
    IO_{max} =
    T_{total} \cdot
    E\bigl[io \mid S_{wide}\bigr].
\]
Такой выбор согласован с операциями объединения и разделения схем: при
объединении ширина полосы растёт, избыточность уменьшается, но вероятность
затронуть схему отказами и ожидаемая стоимость чтения увеличиваются. При
разделении происходит обратный обмен.

Второй глобальной метрикой является суммарная избыточность кластера. Для каждой
схемы \(i\) определим избыточность как разность между физическим объёмом хранения и
объёмом полезных данных:
\[
    Overhead_{sum} = \sum_i \bigl(Space(S_i) - Data_i\bigr),
\]
где \(Data_i\) --- объём данных схемы \(i\).

Для нормировки пространственного давления используются физические ограничения
кластера. Нижняя граница избыточности соответствует хранению всех данных в наиболее
компактной схеме \(\mathrm{LRCC}(d_{max}, 1, 2)\):
\[
    Overhead_{min} = Data_{total} \cdot \frac{3}{d_{max}},
\]
где \(Data_{total} = \sum_i Data_i\) --- суммарный объём полезных данных, а
коэффициент \(3/d_{max}\) отражает долю избыточности при \(l=1, g=2\).
Верхняя граница определяется ёмкостью кластера:
\[
    Overhead_{max} = C_{cluster} - Data_{total},
\]
где \(C_{cluster}\) --- суммарная ёмкость всех дисков кластера.

Нормированные давления по месту и чтению задаются формулами
\[
    P_{space} =
    \frac{Overhead_{sum} - Overhead_{min}}{Overhead_{max} - Overhead_{min}},
    \qquad
    P_{io} =
    \frac{IO_{sum} - IO_{min}}{IO_{max} - IO_{min}}.
\]
Если знаменатель равен нулю, соответствующее давление считается равным нулю.

Итоговая функция состояния имеет вид
\[
    Q = \frac{1}{2}\left(P_{space}^2 + P_{io}^2\right).
\]
Переход считается полезным, если он уменьшает \(Q\). Чтобы учитывать стоимость
фоновой работы, полезность перехода нормируется на объём работы:
\[
    \operatorname{score}(T) =
    \frac{
        Q_{before} - Q_{after}(T)
    }{
        \operatorname{work}(T)
    }.
\]
Величина \(\operatorname{work}(T)\) является скалярной оценкой работы
перекодирования: сколько данных нужно прочитать, записать и удалить. Для
сравнения переходов используется симметричная оценка стоимости:
\(\operatorname{work}(T)=work_{forward}(T)+work_{reverse}(T)\). Фактически
исполняется только выбранное направление, но такая нормировка не поощряет
переходы, которые дешевы при охлаждении и существенно дороже при последующем
нагреве. Критерий не требует вводить произвольные веса между местом, задержкой,
сетью и процессорной работой.

\pagebreak
